%% kapa: exam=1 round=29 year=2018 period=1 form=A subject=civil_law subject_ko=민법 seq=38 number=38
%% source: Q-Net 감정평가사 기출문제 (cst00309) artlSeq=5202808 file=extracted/1008/5202808/2018년 제29회 감정평가사 1차시험 기출문제/2018년 제29회 감정평가사 1차시험 기출문제/1교시 A형.pdf
%% conversion: kordoc 4.24.51 -> tools/gen1.py (auto); review=proofread-sonnet; flags=none
\begin{question}{38}
지상권에 관한 설명으로 옳지 않은 것은? (다툼이 있으면 판례에 따름)
\begin{choices}
  \item 지상권자는 그 권리의 존속기간 내에서 그 토지를 임대할 수 있다.
  \item 지상권이 소멸한 경우 특별한 사정이 없는 한, 지상권자는 건물 기타 공작물이나 수목을 수거하여 토지를 원상에 회복하여야 한다.
  \item 지상권자의 지료지급이 토지소유권의 양도 전후에 걸쳐 2년 이상 연체된 경우, 토지양수인에 대한 연체기간이 2년이 되지 않더라도 토지양수인은 지상권소멸청구를 할 수 있다.
  \item 나대지에 저당권을 설정한 당사자들이 그 목적 토지상에 저당권자 앞으로 저당 토지의 담보가치 저감을 막기 위하여 지상권도 설정한 경우, 저당권의 피담보채권이 시효로 소멸하면 지상권도 소멸한다.
  \item 토지와 그 지상건물이 함께 양도되었다가 채권자취소권의 행사로 그 중 건물에 대해서만 양도가 취소되어 수익자 명의의 소유권이전등기가 말소된 경우, 채무자에게 관습상 법정지상권은 인정되지 않는다.
\end{choices}
\end{question}
\answer{3}
